\sheethead{03 磁耦合元件与变压器}{重点检查：线圈弧形、互感弧线、同名端、变压器左右对称性与间距，以及电压标注相对引脚的位置。}
\begin{center}
\begin{circuitikz}[x=1cm,y=1cm,wire,european resistors,american inductors]
\node[font=\small\bfseries] at (6.2,11.35) {耦合电感};
% 左侧线圈，弧形向外
\draw (4.3,9.8)--(4.3,8.6);
\draw (4.3,8.6)
 .. controls (3.45,8.6) and (3.45,7.8) .. (4.3,7.8)
 .. controls (3.45,7.8) and (3.45,7.0) .. (4.3,7.0)
 .. controls (3.45,7.0) and (3.45,6.2) .. (4.3,6.2)
 .. controls (3.45,6.2) and (3.45,5.4) .. (4.3,5.4)
 -- (4.3,4.3);
\draw (4.3,4.3)--(4.3,3.45);
% 右侧线圈，弧形向外
\draw (8.1,9.8)--(8.1,8.6);
\draw (8.1,8.6)
 .. controls (8.95,8.6) and (8.95,7.8) .. (8.1,7.8)
 .. controls (8.95,7.8) and (8.95,7.0) .. (8.1,7.0)
 .. controls (8.95,7.0) and (8.95,6.2) .. (8.1,6.2)
 .. controls (8.95,6.2) and (8.95,5.4) .. (8.1,5.4)
 -- (8.1,4.3);
\draw (8.1,4.3)--(8.1,3.45);
\fill (4.3,8.52) circle (2.2pt);
\fill (8.1,8.52) circle (2.2pt);
% 单条连续互感弧线
\draw (4.35,8.52) .. controls (5.0,9.75) and (7.4,9.75) .. (8.05,8.52);
\node[font=\normalsize] at (6.2,10.05) {$M$};
\node[font=\small] at (4.3,2.95) {$L_1$};
\node[font=\small] at (8.1,2.95) {$L_2$};
\node[font=\footnotesize,align=center,text width=8cm] at (6.2,1.95) {个人偏好候选：弧形线圈、单条连续的互感弧线、同名端点清楚；线圈间距适中。};

\draw[gray!60,line width=.4pt] (13.3,1.1)--(13.3,11.5);
\node[font=\small\bfseries] at (20.4,11.35) {变压器};
% 原边线圈
\draw (19.6,9.8)--(19.6,8.6);
\draw (19.6,8.6)
 .. controls (18.85,8.6) and (18.85,7.8) .. (19.6,7.8)
 .. controls (18.85,7.8) and (18.85,7.0) .. (19.6,7.0)
 .. controls (18.85,7.0) and (18.85,6.2) .. (19.6,6.2)
 .. controls (18.85,6.2) and (18.85,5.4) .. (19.6,5.4)
 -- (19.6,4.3);
\draw (19.6,4.3)--(19.6,3.45);
% 副边线圈
\draw (21.2,9.8)--(21.2,8.6);
\draw (21.2,8.6)
 .. controls (21.95,8.6) and (21.95,7.8) .. (21.2,7.8)
 .. controls (21.95,7.8) and (21.95,7.0) .. (21.2,7.0)
 .. controls (21.95,7.0) and (21.95,6.2) .. (21.2,6.2)
 .. controls (21.95,6.2) and (21.95,5.4) .. (21.2,5.4)
 -- (21.2,4.3);
\draw (21.2,4.3)--(21.2,3.45);
% 铁芯位于中心，靠近两侧绕组
\draw (20.28,4.55)--(20.28,8.35);
\draw (20.52,4.55)--(20.52,8.35);
\fill (19.6,8.52) circle (2.2pt);
\fill (21.2,8.52) circle (2.2pt);
\node[font=\normalsize] at (20.4,10.25) {$N_1:N_2$};
\node[font=\small] at (17.25,6.6) {$u_1$};
\node[font=\small] at (23.55,6.6) {$u_2$};
\node[font=\small] at (18.1,7.55) {$+$};
\node[font=\small] at (18.1,5.2) {$-$};
\node[font=\small] at (22.7,7.55) {$+$};
\node[font=\small] at (22.7,5.2) {$-$};
\node[font=\footnotesize,align=center,text width=9cm] at (20.4,1.95) {个人偏好候选：原副边关于铁芯近似对称，收紧空白；$u_1,u_2$ 靠各自绕组，$N_1:N_2$ 居中。};
\end{circuitikz}
\end{center}
